\subsection{Encoders: cross-frequency structure, or none}
\label{sec:encoder}

The grid $(C \times 2n_b \times P \times d)$ can be processed by an encoder
that supplies cross-frequency relations itself, or by one that does not.
We build both, with the same tokenizer, the same recipe and the same
mean-pooled linear head, because the difference between them is the
experiment of Section~\ref{sec:content-structure}.

\paragraph{Tri-axial encoder (structure).}
$N$ pre-norm blocks each apply self-attention along one axis at a time ---
time (with rotary position encoding), electrodes (with a learned
per-electrode embedding), and the $2n_b$ frequency rows (with an embedding
of each band's learned center frequency and bandwidth) --- followed by a
feed-forward layer, each as a residual. The frequency sub-layer is plain
attention over rows: it gives the encoder a path along which it can
compute any cross-band relation from the waveform rows on its own, so the
interaction token competes with structure that can re-derive it. We use
$n_b{=}8$, $d{=}128$, $N{=}6$, four heads, 1.6M parameters.

\paragraph{Axis-free encoder (content only).}
The final model removes the frequency sub-layer entirely and folds the band
rows into the spatial one: in each block, attention runs first along time
per (electrode, row), then jointly over the $C \cdot 2n_b$ tokens of each
patch, then through the feed-forward layer. Rows of different bands meet
only inside this folded spatial attention, which carries no band
embedding and no frequency-specific parameters. The encoder therefore has
no pathway of its own for cross-frequency relations; the only place such a
relation can enter is the interaction token of Eq.~(\ref{eq:rotation}),
and switching that token off (Proposition~\ref{prop:recover}) leaves a
model that sees eight bands as eight unrelated channels. Width is raised
to $d{=}192$ to spend the parameters the frequency sub-layer no longer
uses; $N{=}6$, four heads, 2.7M parameters. Section~\ref{sec:design}
reports the alternatives (no folding, $d{=}256$, an explicit
coupling-strength feature) that were tried and not adopted.

\paragraph{The token inside a published encoder.}
To test the same distinction in an encoder we did not design, we add the
tokenizer's rows to CBraMod \citep{wang2025cbramod} without removing
anything: CBraMod's patch embedding (a convolutional stem plus an rFFT
magnitude branch), its criss-cross encoder, positional convolution and
classifier are kept verbatim, and each electrode receives $n_b$ extra
channels carrying either the $n_b$ interaction tokens or, as the matched
control, the $n_b$ waveform tokens of the same frontend. The two arms have
the same parameter count and the same token count and differ only in what
the extra rows carry. CBraMod's own spectral branch remains in place, so
this encoder, unlike the axis-free one, is never starved of
cross-frequency information.
